\documentclass[11pt,a4paper]{article}
\usepackage{fontspec}
\IfFontExistsTF{DejaVu Sans}{\setmainfont{DejaVu Sans}}{\IfFontExistsTF{Segoe UI}{\setmainfont{Segoe UI}}{\setmainfont{Latin Modern Roman}}}
\usepackage[margin=24mm]{geometry}
\usepackage{xcolor}
\usepackage{hyperref}
\usepackage{xurl}
\usepackage{parskip}
\definecolor{kblue}{HTML}{215D8E}
\definecolor{ksoft}{HTML}{EDF3F7}
\newcommand{\spekSlot}[2]{\par\medskip{\small\bfseries\color{kblue}#1}\par\smallskip #2\par}
\newcommand{\spekFeatured}[2]{\par\medskip\noindent\fcolorbox{kblue}{ksoft}{\parbox{\dimexpr\linewidth-2\fboxsep-2\fboxrule}{\textbf{#1}\par\medskip #2}}\par\bigskip}
\hypersetup{colorlinks=true,linkcolor=kblue}
\title{3e année — Étape pédagogique 2}
\date{}
\begin{document}
\maketitle
\noindent\fbox{\parbox{0.95\textwidth}{\textbf{DÉMONSTRATION — NON VALIDÉE SCIENTIFIQUEMENT}}}

\begin{abstract}
Huit scénarios proposés et une semaine de consolidation. Les correspondances avec les cours Ekole restent des candidats à la revue pédagogique; aucune maîtrise ni affectation officielle n’est inférée.
\end{abstract}
\tableofcontents
\newpage
\section{Semaine 10 — La commande de la boulangerie fictive}
\textit{Parcours de notions}

% spek-section-id: week-10

\spekFeatured{Objectif}{Calculer sur des nombres adaptés, représenter la démarche, puis vérifier le résultat en situation.\par{\scriptsize\ttfamily u.w10.objective · objective · design\_proposal · refs: e.w10.plan, e.w10.overlay, e.w10.micro · source-ref: kontexte:grade3/overlay:qc3-commande-de-pains}}
\noindent\begin{minipage}[t]{.48\linewidth}\spekSlot{Préalables}{Préalables précis à définir par l’enseignant; aucune compétence acquise n’est présumée.\par{\scriptsize\ttfamily u.w10.pre · prerequisite · limitation · source-ref: editorial:unresolved-prerequisite}}\end{minipage}\hfill\begin{minipage}[t]{.48\linewidth}\spekSlot{Question de réflexion}{Trace proposée : Calcul, unité et vérification du résultat.\par{\scriptsize\ttfamily u.w10.reflection · reflection · design\_proposal · refs: e.w10.plan, e.w10.overlay, e.w10.micro · source-ref: kontexte:grade3/overlay:qc3-commande-de-pains}}\end{minipage}\par\bigskip
\spekSlot{Étapes}{Les commandes reçues dépassent les quantités préparées.\par{\scriptsize\ttfamily u.w10.situation · step · method · refs: e.w10.plan, e.w10.overlay, e.w10.micro · source-ref: kontexte:grade3/overlay:qc3-commande-de-pains}\par\smallskip Livrable proposé : bon de commande et calcul des lots\par{\scriptsize\ttfamily u.w10.deliverable · step · design\_proposal · refs: e.w10.plan, e.w10.overlay, e.w10.micro · source-ref: kontexte:grade3/overlay:qc3-commande-de-pains}\par\smallskip Reformuler les consignes et produire un écrit ou une intervention orale pour résoudre le problème.\par{\scriptsize\ttfamily u.w10.action1 · step · method · refs: e.w10.plan, e.w10.overlay, e.w10.micro · source-ref: kontexte:grade3/overlay:qc3-commande-de-pains}\par\smallskip Formuler une hypothèse, comparer des observations et consigner des données.\par{\scriptsize\ttfamily u.w10.action2 · step · method · refs: e.w10.plan, e.w10.overlay, e.w10.micro · source-ref: kontexte:grade3/overlay:qc3-commande-de-pains}}
\spekSlot{Limites}{Proposition locale à revoir : cette fiche ne valide ni programme officiel, ni couverture réelle, ni maîtrise d’un élève.\par{\scriptsize\ttfamily u.w10.limits · qualification · limitation · source-ref: editorial:review-required}}
\section{Semaine 11 — Deux bacs de haricots}
\textit{Parcours de notions}

% spek-section-id: week-11

\spekFeatured{Objectif}{Formuler une hypothèse, comparer des observations et consigner des données.\par{\scriptsize\ttfamily u.w11.objective · objective · design\_proposal · refs: e.w11.plan, e.w11.overlay, e.w11.micro · source-ref: kontexte:grade3/overlay:qc3-semis-compare}}
\noindent\begin{minipage}[t]{.48\linewidth}\spekSlot{Préalables}{Préalables précis à définir par l’enseignant; aucune compétence acquise n’est présumée.\par{\scriptsize\ttfamily u.w11.pre · prerequisite · limitation · source-ref: editorial:unresolved-prerequisite}}\end{minipage}\hfill\begin{minipage}[t]{.48\linewidth}\spekSlot{Question de réflexion}{Trace proposée : Journal daté ou tableau de mesures.\par{\scriptsize\ttfamily u.w11.reflection · reflection · design\_proposal · refs: e.w11.plan, e.w11.overlay, e.w11.micro · source-ref: kontexte:grade3/overlay:qc3-semis-compare}}\end{minipage}\par\bigskip
\spekSlot{Étapes}{Les semis présentent des croissances différentes.\par{\scriptsize\ttfamily u.w11.situation · step · method · refs: e.w11.plan, e.w11.overlay, e.w11.micro · source-ref: kontexte:grade3/overlay:qc3-semis-compare}\par\smallskip Livrable proposé : tableau de mesures et hypothèse\par{\scriptsize\ttfamily u.w11.deliverable · step · design\_proposal · refs: e.w11.plan, e.w11.overlay, e.w11.micro · source-ref: kontexte:grade3/overlay:qc3-semis-compare}\par\smallskip Calculer sur des nombres adaptés, représenter la démarche, puis vérifier le résultat en situation.\par{\scriptsize\ttfamily u.w11.action1 · step · method · refs: e.w11.plan, e.w11.overlay, e.w11.micro · source-ref: kontexte:grade3/overlay:qc3-semis-compare}}
\spekSlot{Limites}{Proposition locale à revoir : cette fiche ne valide ni programme officiel, ni couverture réelle, ni maîtrise d’un élève.\par{\scriptsize\ttfamily u.w11.limits · qualification · limitation · source-ref: editorial:review-required}}
\section{Semaine 12 — Bienvenue à notre école}
\textit{Parcours de notions}

% spek-section-id: week-12

\spekFeatured{Objectif}{Reformuler les consignes et produire un écrit ou une intervention orale pour résoudre le problème.\par{\scriptsize\ttfamily u.w12.objective · objective · design\_proposal · refs: e.w12.plan, e.w12.overlay, e.w12.micro · source-ref: kontexte:grade3/overlay:qc3-accueil-nouveaux}}
\noindent\begin{minipage}[t]{.48\linewidth}\spekSlot{Préalables}{Préalables précis à définir par l’enseignant; aucune compétence acquise n’est présumée.\par{\scriptsize\ttfamily u.w12.pre · prerequisite · limitation · source-ref: editorial:unresolved-prerequisite}}\end{minipage}\hfill\begin{minipage}[t]{.48\linewidth}\spekSlot{Question de réflexion}{Trace proposée : Une version annotée du texte et une version révisée.\par{\scriptsize\ttfamily u.w12.reflection · reflection · design\_proposal · refs: e.w12.plan, e.w12.overlay, e.w12.micro · source-ref: kontexte:grade3/overlay:qc3-accueil-nouveaux}}\end{minipage}\par\bigskip
\spekSlot{Étapes}{Des renseignements nécessaires aux nouveaux élèves sont dispersés.\par{\scriptsize\ttfamily u.w12.situation · step · method · refs: e.w12.plan, e.w12.overlay, e.w12.micro · source-ref: kontexte:grade3/overlay:qc3-accueil-nouveaux}\par\smallskip Livrable proposé : guide scolaire illustré\par{\scriptsize\ttfamily u.w12.deliverable · step · design\_proposal · refs: e.w12.plan, e.w12.overlay, e.w12.micro · source-ref: kontexte:grade3/overlay:qc3-accueil-nouveaux}\par\smallskip Lire ou écouter une très courte formulation anglaise utile à un destinataire réel ou simulé.\par{\scriptsize\ttfamily u.w12.action1 · step · method · refs: e.w12.plan, e.w12.overlay, e.w12.micro · source-ref: kontexte:grade3/overlay:qc3-accueil-nouveaux}\par\smallskip Localiser les lieux ou documents dans un contexte québécois, distinguer les sources.\par{\scriptsize\ttfamily u.w12.action2 · step · method · refs: e.w12.plan, e.w12.overlay, e.w12.micro · source-ref: kontexte:grade3/overlay:qc3-accueil-nouveaux}}
\spekSlot{Limites}{Proposition locale à revoir : cette fiche ne valide ni programme officiel, ni couverture réelle, ni maîtrise d’un élève.\par{\scriptsize\ttfamily u.w12.limits · qualification · limitation · source-ref: editorial:review-required}}
\section{Semaine 13 — L’invitation fr-CA / en-CA}
\textit{Parcours de notions}

% spek-section-id: week-13

\spekFeatured{Objectif}{Reformuler les consignes et produire un écrit ou une intervention orale pour résoudre le problème.\par{\scriptsize\ttfamily u.w13.objective · objective · design\_proposal · refs: e.w13.plan, e.w13.overlay, e.w13.micro · source-ref: kontexte:grade3/overlay:qc3-invitation-multilingue}}
\noindent\begin{minipage}[t]{.48\linewidth}\spekSlot{Préalables}{Préalables précis à définir par l’enseignant; aucune compétence acquise n’est présumée.\par{\scriptsize\ttfamily u.w13.pre · prerequisite · limitation · source-ref: editorial:unresolved-prerequisite}}\end{minipage}\hfill\begin{minipage}[t]{.48\linewidth}\spekSlot{Question de réflexion}{Trace proposée : Une version annotée du texte et une version révisée.\par{\scriptsize\ttfamily u.w13.reflection · reflection · design\_proposal · refs: e.w13.plan, e.w13.overlay, e.w13.micro · source-ref: kontexte:grade3/overlay:qc3-invitation-multilingue}}\end{minipage}\par\bigskip
\spekSlot{Étapes}{La version anglaise omet une donnée importante.\par{\scriptsize\ttfamily u.w13.situation · step · method · refs: e.w13.plan, e.w13.overlay, e.w13.micro · source-ref: kontexte:grade3/overlay:qc3-invitation-multilingue}\par\smallskip Livrable proposé : invitation bilingue validée\par{\scriptsize\ttfamily u.w13.deliverable · step · design\_proposal · refs: e.w13.plan, e.w13.overlay, e.w13.micro · source-ref: kontexte:grade3/overlay:qc3-invitation-multilingue}\par\smallskip Lire ou écouter une très courte formulation anglaise utile à un destinataire réel ou simulé.\par{\scriptsize\ttfamily u.w13.action1 · step · method · refs: e.w13.plan, e.w13.overlay, e.w13.micro · source-ref: kontexte:grade3/overlay:qc3-invitation-multilingue}\par\smallskip Réaliser une production plastique et expliquer les choix de matériaux.\par{\scriptsize\ttfamily u.w13.action2 · step · method · refs: e.w13.plan, e.w13.overlay, e.w13.micro · source-ref: kontexte:grade3/overlay:qc3-invitation-multilingue}}
\spekSlot{Limites}{Proposition locale à revoir : cette fiche ne valide ni programme officiel, ni couverture réelle, ni maîtrise d’un élève.\par{\scriptsize\ttfamily u.w13.limits · qualification · limitation · source-ref: editorial:review-required}}
\section{Semaine 14 — Le budget des kiosques}
\textit{Parcours de notions}

% spek-section-id: week-14

\spekFeatured{Objectif}{Calculer sur des nombres adaptés, représenter la démarche, puis vérifier le résultat en situation.\par{\scriptsize\ttfamily u.w14.objective · objective · design\_proposal · refs: e.w14.plan, e.w14.overlay, e.w14.micro · source-ref: kontexte:grade3/overlay:qc3-budget-marche}}
\noindent\begin{minipage}[t]{.48\linewidth}\spekSlot{Préalables}{Préalables précis à définir par l’enseignant; aucune compétence acquise n’est présumée.\par{\scriptsize\ttfamily u.w14.pre · prerequisite · limitation · source-ref: editorial:unresolved-prerequisite}}\end{minipage}\hfill\begin{minipage}[t]{.48\linewidth}\spekSlot{Question de réflexion}{Trace proposée : Calcul, unité et vérification du résultat.\par{\scriptsize\ttfamily u.w14.reflection · reflection · design\_proposal · refs: e.w14.plan, e.w14.overlay, e.w14.micro · source-ref: kontexte:grade3/overlay:qc3-budget-marche}}\end{minipage}\par\bigskip
\spekSlot{Étapes}{Le budget de l’équipe est inférieur aux achats souhaités.\par{\scriptsize\ttfamily u.w14.situation · step · method · refs: e.w14.plan, e.w14.overlay, e.w14.micro · source-ref: kontexte:grade3/overlay:qc3-budget-marche}\par\smallskip Livrable proposé : budget avec choix justifiés\par{\scriptsize\ttfamily u.w14.deliverable · step · design\_proposal · refs: e.w14.plan, e.w14.overlay, e.w14.micro · source-ref: kontexte:grade3/overlay:qc3-budget-marche}\par\smallskip Reformuler les consignes et produire un écrit ou une intervention orale pour résoudre le problème.\par{\scriptsize\ttfamily u.w14.action1 · step · method · refs: e.w14.plan, e.w14.overlay, e.w14.micro · source-ref: kontexte:grade3/overlay:qc3-budget-marche}}
\spekSlot{Limites}{Proposition locale à revoir : cette fiche ne valide ni programme officiel, ni couverture réelle, ni maîtrise d’un élève.\par{\scriptsize\ttfamily u.w14.limits · qualification · limitation · source-ref: editorial:review-required}}
\section{Semaine 15 — Le potager des quatre saisons}
\textit{Parcours de notions}

% spek-section-id: week-15

\spekFeatured{Objectif}{Formuler une hypothèse, comparer des observations et consigner des données.\par{\scriptsize\ttfamily u.w15.objective · objective · design\_proposal · refs: e.w15.plan, e.w15.overlay, e.w15.micro · source-ref: kontexte:grade3/overlay:qc3-jardin-saisonnier}}
\noindent\begin{minipage}[t]{.48\linewidth}\spekSlot{Préalables}{Préalables précis à définir par l’enseignant; aucune compétence acquise n’est présumée.\par{\scriptsize\ttfamily u.w15.pre · prerequisite · limitation · source-ref: editorial:unresolved-prerequisite}}\end{minipage}\hfill\begin{minipage}[t]{.48\linewidth}\spekSlot{Question de réflexion}{Trace proposée : Journal daté ou tableau de mesures.\par{\scriptsize\ttfamily u.w15.reflection · reflection · design\_proposal · refs: e.w15.plan, e.w15.overlay, e.w15.micro · source-ref: kontexte:grade3/overlay:qc3-jardin-saisonnier}}\end{minipage}\par\bigskip
\spekSlot{Étapes}{Il manque des renseignements sur arrosage et luminosité.\par{\scriptsize\ttfamily u.w15.situation · step · method · refs: e.w15.plan, e.w15.overlay, e.w15.micro · source-ref: kontexte:grade3/overlay:qc3-jardin-saisonnier}\par\smallskip Livrable proposé : journal de jardinage\par{\scriptsize\ttfamily u.w15.deliverable · step · design\_proposal · refs: e.w15.plan, e.w15.overlay, e.w15.micro · source-ref: kontexte:grade3/overlay:qc3-jardin-saisonnier}\par\smallskip Calculer sur des nombres adaptés, représenter la démarche, puis vérifier le résultat en situation.\par{\scriptsize\ttfamily u.w15.action1 · step · method · refs: e.w15.plan, e.w15.overlay, e.w15.micro · source-ref: kontexte:grade3/overlay:qc3-jardin-saisonnier}\par\smallskip Reformuler les consignes et produire un écrit ou une intervention orale pour résoudre le problème.\par{\scriptsize\ttfamily u.w15.action2 · step · method · refs: e.w15.plan, e.w15.overlay, e.w15.micro · source-ref: kontexte:grade3/overlay:qc3-jardin-saisonnier}}
\spekSlot{Limites}{Proposition locale à revoir : cette fiche ne valide ni programme officiel, ni couverture réelle, ni maîtrise d’un élève.\par{\scriptsize\ttfamily u.w15.limits · qualification · limitation · source-ref: editorial:review-required}}
\section{Semaine 16 — Les matériaux partagés}
\textit{Parcours de notions}

% spek-section-id: week-16

\spekFeatured{Objectif}{Calculer sur des nombres adaptés, représenter la démarche, puis vérifier le résultat en situation.\par{\scriptsize\ttfamily u.w16.objective · objective · design\_proposal · refs: e.w16.plan, e.w16.overlay, e.w16.micro · source-ref: kontexte:grade3/overlay:qc3-ressources-a-partager}}
\noindent\begin{minipage}[t]{.48\linewidth}\spekSlot{Préalables}{Préalables précis à définir par l’enseignant; aucune compétence acquise n’est présumée.\par{\scriptsize\ttfamily u.w16.pre · prerequisite · limitation · source-ref: editorial:unresolved-prerequisite}}\end{minipage}\hfill\begin{minipage}[t]{.48\linewidth}\spekSlot{Question de réflexion}{Trace proposée : Calcul, unité et vérification du résultat.\par{\scriptsize\ttfamily u.w16.reflection · reflection · design\_proposal · refs: e.w16.plan, e.w16.overlay, e.w16.micro · source-ref: kontexte:grade3/overlay:qc3-ressources-a-partager}}\end{minipage}\par\bigskip
\spekSlot{Étapes}{Les stocks ne suffisent pas à toutes les demandes.\par{\scriptsize\ttfamily u.w16.situation · step · method · refs: e.w16.plan, e.w16.overlay, e.w16.micro · source-ref: kontexte:grade3/overlay:qc3-ressources-a-partager}\par\smallskip Livrable proposé : tableau de distribution accepté\par{\scriptsize\ttfamily u.w16.deliverable · step · design\_proposal · refs: e.w16.plan, e.w16.overlay, e.w16.micro · source-ref: kontexte:grade3/overlay:qc3-ressources-a-partager}\par\smallskip Comparer des points de vue, expliquer une décision et écouter les autres sans les catégoriser.\par{\scriptsize\ttfamily u.w16.action1 · step · method · refs: e.w16.plan, e.w16.overlay, e.w16.micro · source-ref: kontexte:grade3/overlay:qc3-ressources-a-partager}}
\spekSlot{Limites}{Proposition locale à revoir : cette fiche ne valide ni programme officiel, ni couverture réelle, ni maîtrise d’un élève.\par{\scriptsize\ttfamily u.w16.limits · qualification · limitation · source-ref: editorial:review-required}}
\section{Semaine 17 — L’affiche du service d’entraide}
\textit{Parcours de notions}

% spek-section-id: week-17

\spekFeatured{Objectif}{Reformuler les consignes et produire un écrit ou une intervention orale pour résoudre le problème.\par{\scriptsize\ttfamily u.w17.objective · objective · design\_proposal · refs: e.w17.plan, e.w17.overlay, e.w17.micro · source-ref: kontexte:grade3/overlay:qc3-affiche-service}}
\noindent\begin{minipage}[t]{.48\linewidth}\spekSlot{Préalables}{Préalables précis à définir par l’enseignant; aucune compétence acquise n’est présumée.\par{\scriptsize\ttfamily u.w17.pre · prerequisite · limitation · source-ref: editorial:unresolved-prerequisite}}\end{minipage}\hfill\begin{minipage}[t]{.48\linewidth}\spekSlot{Question de réflexion}{Trace proposée : Une version annotée du texte et une version révisée.\par{\scriptsize\ttfamily u.w17.reflection · reflection · design\_proposal · refs: e.w17.plan, e.w17.overlay, e.w17.micro · source-ref: kontexte:grade3/overlay:qc3-affiche-service}}\end{minipage}\par\bigskip
\spekSlot{Étapes}{L’affiche actuelle n’explique pas comment demander de l’aide.\par{\scriptsize\ttfamily u.w17.situation · step · method · refs: e.w17.plan, e.w17.overlay, e.w17.micro · source-ref: kontexte:grade3/overlay:qc3-affiche-service}\par\smallskip Livrable proposé : affiche accessible et bilingue simple\par{\scriptsize\ttfamily u.w17.deliverable · step · design\_proposal · refs: e.w17.plan, e.w17.overlay, e.w17.micro · source-ref: kontexte:grade3/overlay:qc3-affiche-service}\par\smallskip Lire ou écouter une très courte formulation anglaise utile à un destinataire réel ou simulé.\par{\scriptsize\ttfamily u.w17.action1 · step · method · refs: e.w17.plan, e.w17.overlay, e.w17.micro · source-ref: kontexte:grade3/overlay:qc3-affiche-service}\par\smallskip Réaliser une production plastique et expliquer les choix de matériaux.\par{\scriptsize\ttfamily u.w17.action2 · step · method · refs: e.w17.plan, e.w17.overlay, e.w17.micro · source-ref: kontexte:grade3/overlay:qc3-affiche-service}}
\spekSlot{Limites}{Proposition locale à revoir : cette fiche ne valide ni programme officiel, ni couverture réelle, ni maîtrise d’un élève.\par{\scriptsize\ttfamily u.w17.limits · qualification · limitation · source-ref: editorial:review-required}}
\section{Semaine 18 — Consolidation et reprises}
\textit{Parcours de notions}

% spek-section-id: week-18

\spekFeatured{Objectif}{consolidation, reprises et évaluation sans nouveau scénario obligatoire\par{\scriptsize\ttfamily u.w18.objective · objective · design\_proposal · refs: e.w18.plan, e.w18.micro · source-ref: kontexte:grade3/week:18}}
\noindent\begin{minipage}[t]{.48\linewidth}\spekSlot{Préalables}{Préalables précis à définir par l’enseignant; aucune compétence acquise n’est présumée.\par{\scriptsize\ttfamily u.w18.pre · prerequisite · limitation · source-ref: editorial:unresolved-prerequisite}}\end{minipage}\hfill\begin{minipage}[t]{.48\linewidth}\spekSlot{Question de réflexion}{Trace proposée : justification d’accord simple\par{\scriptsize\ttfamily u.w18.reflection · reflection · design\_proposal · refs: e.w18.plan, e.w18.micro · source-ref: kontexte:grade3/week:18}}\end{minipage}\par\bigskip
\spekSlot{Étapes}{consolidation, reprises et évaluation sans nouveau scénario obligatoire\par{\scriptsize\ttfamily u.w18.consolidation · step · method · refs: e.w18.plan, e.w18.micro · source-ref: kontexte:grade3/week:18}\par\smallskip Faire annoter l’élément linguistique ciblé et produire une version révisée.\par{\scriptsize\ttfamily u.w18.micro · step · design\_proposal · refs: e.w18.plan, e.w18.micro · source-ref: kontexte:grade3/week:18}}
\spekSlot{Limites}{Proposition locale à revoir : cette fiche ne valide ni programme officiel, ni couverture réelle, ni maîtrise d’un élève.\par{\scriptsize\ttfamily u.w18.limits · qualification · limitation · source-ref: editorial:review-required}}
\section*{Traçabilité}
\small Statut: Identité canonique non vérifiée
\par Origine: \nolinkurl{sha256:032ea6cad0d27eec5ded58622005dc5ceb55b78b4d6fbafe10d39fcc9f74d364}
\par Réalisations: \nolinkurl{sha256:f9278044e6c8c3642bedd3b8fd2390fc3ece5f2b6666ca3b7eeff23571350de5}
\par Mises en page: \nolinkurl{sha256:77c5d1a2639afe1a346253f5ddc5f5723f6e2a25931f714bb7cdcd647b9319eb}
\end{document}
